% =====================================================================
% 10_perspectives.tex — perspectives et prolongements
% =====================================================================
\chapter{Perspectives}
\label{chap:perspectives}

\section{Introduction}

Les perspectives ne corrigent pas des erreurs : elles prolongent un projet qui
atteint son périmètre. Elles se répartissent en quatre axes : élargir le modèle
métier, renforcer la sécurité et l'exploitation, ajouter une interface et
mettre en place des tests systématiques. Chaque proposition indique ce qu'elle
apporterait et ce qu'elle demanderait comme travail.

\section{Élargir le modèle métier}

\subsection{Opérations polymorphes sur tous les types de comptes}

Aujourd'hui, \code{operation\_type.ref\_compte} est un
\code{REF cptCourant\_type}. Cette décision rend l'interdiction structurelle,
mais elle empêche les dépôts sur un compte d'épargne et les virements vers un
compte à terme. Une évolution consisterait à créer une super-classe
\code{operation\_type} non instanciable, avec des sous-types
\code{operation\_courant\_type}, \code{operation\_epargne\_type} et
\code{operation\_terme\_type}. Une autre solution serait une table d'opérations
avec un discriminant et une référence générique vers \code{compte\_type}.

Cette évolution demanderait de revoir les triggers, les méthodes de calcul de
solde et les requêtes d'agrégation. Elle apporterait une couverture complète
des mouvements d'épargne et de terme, ainsi qu'un vrai polymorphisme des
opérations.

\subsection{Virement entre comptes de nature différente}

Le virement actuel relie deux comptes courants. Un virement courant vers
épargne pourrait être autorisé si le compte destination accepte un crédit et si
le plafond de dépôt est respecté. Il faudrait alors vérifier la nature des
comptes, appliquer les règles de plafond et adapter la commission. La
commission inter-agences resterait calculée sur le compte source.

\subsection{Responsable d'agence et intérim}

La méthode \code{mutation} signale qu'un conseiller responsable quitte son
agence, mais ne désigne pas de remplaçant. On pourrait ajouter une table
\code{interim\_agence} ou une procédure \code{nomme\_responsable} qui vérifie
l'appartenance du nouveau responsable à l'agence, la date de prise de fonction
et l'absence de conflit. Cette évolution relèverait d'une règle de management
explicite, validée par le métier.

\subsection{Plafond des dépôts d'épargne}

Le plafond de dépôt est présent comme attribut, mais aucun mouvement d'épargne
n'est créé dans le périmètre actuel. Si les opérations d'épargne étaient
ajoutées, il faudrait un contrôle \code{solde + montant <= plafond\_depot},
soit par contrainte, soit par trigger. La solution dépendrait du niveau de
confiance accordé aux insertions directes.

\subsection{Gestion des décès, successions et comptes joints}

Le modèle ne gère pas les comptes joints, les successions ni les procurations.
On pourrait ajouter une table de titulaires multiples, des rôles datés et des
règles de blocage. Ces ajouts toucheraient la clé primaire des comptes, la
notion de titulaire et les droits de signature. Ils demanderaient une étude
juridique et métier, au-delà du périmètre technique du mini-projet.

\section{Renforcer la sécurité et l'exploitation}

\subsection{Séparation des rôles et privilèges minimaux}

Le projet utilise un seul utilisateur, \code{MASTER1}, avec des privilèges
étendus. Une base réelle séparerait les rôles : administrateur, gestionnaire
d'agence, conseiller, auditeur, application. Chaque rôle recevrait uniquement
les droits nécessaires, et les accès seraient journalisés. Les vues pourraient
masquer les données sensibles, comme le numéro de sécurité sociale.

\subsection{Mots de passe et chiffrement}

Le mot de passe \code{system} est pédagogique. Il faudrait un mot de passe
robuste, une politique d'expiration, et éventuellement le chiffrement des
colonnes sensibles. Oracle propose des mécanismes de chiffrement transparent
et de masquage de données ; leur mise en place demanderait une licence et une
configuration adaptées.

\subsection{Audit autonome}

Les triggers d'audit écrivent dans la même transaction que la modification.
Un \code{ROLLBACK} annule donc la trace. Pour un audit de sécurité, il faudrait
une journalisation autonome, par exemple avec \code{DBMS\_SCHEDULER}, une table
d'audit hors transaction ou Oracle Audit Vault. Cette évolution changerait la
nature de l'audit : il ne serait plus transactionnel mais persistant.

\subsection{Surveillance et alertes}

Les alertes sont stockées dans une table et consultées par requête. On pourrait
ajouter une notification par courriel ou par file d'attente, un tableau de bord
et un acquittement des alertes. Un job planifié pourrait exécuter
\code{scan\_echeances} chaque nuit et marquer les alertes traitées.

\section{Ajouter une interface et des services}

\subsection{Interface graphique}

Le projet n'a pas d'interface. Une application web ou bureau pourrait exposer
les clients, les comptes, les opérations et les alertes. Elle appellerait les
procédures existantes plutôt que d'écrire directement dans les tables, afin de
conserver les contrôles. Les méthodes \code{passeMvt} et \code{virement}
constitueraient la couche de service.

\subsection{API et séparation des couches}

Une API REST pourrait exposer les opérations de lecture et d'écriture. Il
faudrait alors gérer l'authentification, les autorisations, la validation des
entrées et les transactions. La couche base resterait responsable des règles
métier, mais l'API devrait traduire les erreurs applicatives
(\code{ORA-20001} à \code{ORA-20030}) en messages compréhensibles.

\subsection{États et exports}

On pourrait ajouter des états imprimables : relevé de compte, attestation de
solde, historique des opérations, liste des alertes. Ces états utiliseraient
les vues et les collections existantes, avec une mise en forme adaptée. Un
export CSV ou PDF compléterait l'interface.

\section{Tester, mesurer et industrialiser}

\subsection{Tests automatisés}

Les tests actuels sont des scripts SQL commentés. On pourrait les organiser
avec un cadre de test comme \code{utPLSQL}, qui vérifie les résultats attendus
et produit un rapport. Les tests couvriraient les contraintes, les triggers,
les méthodes, les virements et les mutations. Ils seraient exécutés à chaque
modification du schéma.

\subsection{Tests de concurrence}

Le projet ne teste pas deux transactions simultanées. Il faudrait vérifier le
comportement du plafond journalier, des mises à jour de solde et des mutations
lorsque deux sessions travaillent en parallèle. On pourrait utiliser
\code{DBMS\_LOCK}, des niveaux d'isolation adaptés et des tests de verrouillage.
C'est un point important pour une base bancaire réelle.

\subsection{Migrations versionnées}

Le script unique supprime et recrée tout. Pour une exploitation, il faudrait
des migrations versionnées : un historique des changements, des scripts de
montée de version et des scripts de retour arrière. Des outils comme Flyway ou
Liquibase pourraient être utilisés, avec une validation en préproduction.

\subsection{Performances et index}

Le projet possède 21 index, principalement ceux des clés primaires et des
contraintes d'unicité. Une charge réelle demanderait une analyse des plans
d'exécution, des index sur les colonnes de recherche, des statistiques à jour
et peut-être des vues matérialisées pour les agrégats. Les requêtes du rapport
serviraient de base à cette analyse.

\subsection{Documentation et formation}

Le rapport, le guide de soutenance et les commentaires SQL constituent une
première documentation. Une base réelle demanderait un dictionnaire de données,
un manuel d'exploitation, des procédures de sauvegarde et de restauration, et
une formation des utilisateurs. La documentation des hypothèses et des limites
serait maintenue à jour.

\section{Priorisation}

Toutes ces perspectives ne se valent pas. Les plus rentables à court terme
sont les tests automatisés, la séparation des rôles et un interface de
consultation. Les opérations polymorphes et le virement multi-types demandent
une refonte du modèle et des triggers ; elles constituent la prochaine étape
technique naturelle. La gestion des successions et des comptes joints relève
d'une décision métier et juridique plus large.

\begin{table}[htbp]
  \centering
  \small
  \caption{Priorisation indicative des perspectives.}
  \label{tab:priorisation-perspectives}
  \begin{tabularx}{\textwidth}{@{}L{4.4cm}L{2.4cm}X@{}}
    \toprule
    \entete{Perspective} & \entete{Horizon} & \entete{Justification} \\
    \midrule
    Tests automatisés & Court terme &
    Sécurise les modifications et documente les résultats. \\
    Séparation des rôles & Court terme &
    Réduit les risques et prépare une exploitation réelle. \\
    Interface de consultation & Court terme &
    Rend les données accessibles sans SQL. \\
    Opérations polymorphes & Moyen terme &
    Étend le périmètre métier et exploite le modèle objet. \\
    Audit autonome & Moyen terme &
    Conserve les traces malgré les rollbacks. \\
    Tests de concurrence & Moyen terme &
    Vérifie la robustesse des triggers et des soldes. \\
    Migrations versionnées & Moyen terme &
    Prépare une exploitation professionnelle. \\
    Comptes joints et successions & Long terme &
    Dépend de règles métier et juridiques. \\
    \bottomrule
  \end{tabularx}
\end{table}

\section{Conclusion des perspectives}

Le projet est complet dans son périmètre, mais il n'est pas figé. Les
perspectives montrent que les choix actuels peuvent évoluer sans être remis en
cause : le modèle objet facilite l'ajout de sous-types, les procédures
centralisent les règles et les triggers protègent les insertions directes. La
priorité doit aller à la sécurité, aux tests et à l'interface, avant d'élargir
le modèle métier. Cette progression correspond à une industrialisation
raisonnable d'un prototype pédagogique.
